\documentclass[10pt,a4paper]{amsart}
\usepackage[margin=19mm]{geometry}
\usepackage{amsmath,amssymb,mathtools}
\usepackage[scr=rsfs,cal=euler]{mathalfa}
\usepackage{xfrac}
\usepackage[unicode,hidelinks,bookmarks=false]{hyperref}
\newcommand{\ie}{i.e.\ }
\newcommand{\cf}{cf.\ }
\newcommand{\resp}{respectively\ }
\newcommand{\Subsection}{\subsection}
\providecommand{\nobreakdash}{\nobreak\hbox{-}}
\setlength{\emergencystretch}{2em}
\multlinegap0pt
\usepackage{amssymb}
\usepackage{xcolor}
\usepackage{algorithm}\usepackage{algpseudocode}
\newtheorem{lemma}{Lemma}
\newtheorem{proposition}{Proposition}
\newtheorem{corollary}{Corollary}
\newtheorem{theorem}{Theorem}
\title{Tests for the mean of high-dimensional data}
\author{Dietmar Ferger}
\date{Source: arXiv:2605.16033v1 (CC BY 4.0)}
\begin{document}
\maketitle

\noindent\emph{Source and licence.} Tests for the mean of high-dimensional data. Version arXiv:2605.16033v1 (CC BY 4.0), distributed by arXiv under the Creative Commons Attribution 4.0 International licence, http://creativecommons.org/licenses/by/4.0/, as recorded in the arXiv licence metadata for this exact version and shown on the abstract page https://arxiv.org/abs/2605.16033v1. Full licence notice: \texttt{licenses/arxiv-2605.16033-CC-BY-4.0.txt}.\par\smallskip

\noindent\textit{Editorial scope.} Counted unit: the article's theorem bootstraphigh (source0.tex lines 532--538). The article's complete original proof of it is delivered with it (source0.tex lines 540--619, one \texttt{proof} environment of 80 lines). No mathematics is written, repaired or re-derived: the statement and the proof are the article's own lines, in the article's own notation, and no retained line is substituted or re-flowed.  Retained uncounted as the source-local context the proof needs: lines 209--228; lines 244--246; lines 248--250; lines 263--265; lines 273--339; lines 394--396; lines 410--412; lines 452--468; lines 506--511.  Nothing was omitted from any retained span, so the package carries no omission record.  No retained line contains a citation, so the wrapper carries no bibliography and no bibliography entry.  No retained line contains a figure environment, an image-inclusion command or drawing code, so no image is delivered and the package declares no asset dependency.  Editorial formatting only, in addition: an AMS \texttt{amsart} preamble that restates the article's own declarations and macros used by the retained lines, namely the environments \texttt{lemma}, \texttt{proposition}, \texttt{corollary}, \texttt{theorem} on separate counters, together with the standard packages those lines need.  The retained environments print in this document as Lemma 1, Proposition 1, Corollary 1, Theorem 1 in order of appearance, whereas the article prints the same items with the numbering of its own full document.  The complete native source of the article as distributed in the arXiv e-print for this exact version is bundled unchanged as \texttt{sources/200703/source0.tex}.
\begin{lemma} \label{convVn} For every $n \in \mathbb{N}$ let $X_{n1},\ldots,X_{nn}$ be i.i.d. random vectors in $\mathbb{R}^d$ with
$\mathbb{E}[X_{n1}]=0$ and $\mathbb{E}[||X_{n1}||^2]< \infty$. If
\begin{equation}
 \text{Cov}[X_{n1}] \rightarrow \Gamma
\end{equation}
and
\begin{equation} \label{Lindebergconditionwithrationals}
 L_n(\epsilon) := \mathbb{E}[||X_{n1}||^2 1_{\{||X_{n1}||> \sqrt{n} \epsilon\}}] \rightarrow 0 \quad \text{for all positive rational }  \epsilon,
\end{equation}
then $\overline{X}_n := n^{-1} \sum_{i=1}^n X_{ni}$ satisfies the following limit-law:
$$
 V_n:= n ||\overline{X}_n||^2 \stackrel{\mathcal{D}}{\rightarrow} V=\sum_{i=1}^p \lambda_i N_i^2.
$$
Here, $\lambda_1,\dots,\lambda_d$ are the eigenvalues of $\Gamma$ and $N_1,\ldots,N_d$ are i.i.d. standard normal. As we know
this is equivalent to
$$
  \sup_{x \in \mathbb{R}}|\mathbb{P}(V_n \le x)-F(x)| \rightarrow 0, \;  n \rightarrow \infty,
$$
where $F$ is the distribution function of $V$.
\end{lemma}

\begin{equation} \label{defQn}
 Q_n(\omega):= n^{-1} \sum_{i=1}^n \delta_{X_i(\omega)}
\end{equation}

\begin{equation} \label{distribbootstrap}
\mathbb{P}^*\circ (\xi_{n1}(\omega),\ldots,\xi_{nn}(\omega))^{-1}=Q_n(\omega)\otimes \cdots \otimes Q_n(\omega).
\end{equation}

\begin{equation} \label{bootstrapconv}
 V_n^*(\omega)= n ||\overline{\xi}_n(\omega)||^2 \stackrel{\mathcal{D}}{\rightarrow} V=\sum_{i=1}^d \lambda_i N_i^2 \quad \text{for } \mathbb{P}\text{-almost all } \omega \in \Omega.
\end{equation}

\begin{proof} For each fixed $\omega \in \Omega$ the array $X_{ni}(\omega):=\xi_{ni}(\omega)- \overline{X}_n(\omega), 1 \le i \le n, n \in \mathbb{N},$ consists of random variables $X_{ni}(\omega)$, which are defined on $(\Omega^*,\mathcal{A}^*,\mathbb{P}^*)$ and which by (\ref{distribbootstrap}) are rowwise i.i.d. under $\mathbb{P}^*$. In the sequel we check that the array $(X_{ni}(\omega): 1 \le i \le n, n \in \mathbb{N})$ satisfies the remaining conditions of Lemma \ref{convVn} for $\mathbb{P}-$almost all $\omega \in \Omega.$ To start with
$$
   \mathbb{E}^*[\xi_{ni}(\omega)]=\int_{\Omega^*} \xi_{ni}(\omega)(\omega^*) \mathbb{P}^*(d\omega^*)= \int_{\mathbb{R}^p} x  Q_n(\omega)(dx)=n^{-1} \sum_{i=1}^n X_i(\omega) = \overline{X}_n(\omega),
$$
where the second equality holds by Change of Variable plus (\ref{distribbootstrap}) and the third equality follows from the definition (\ref{defQn}) of $Q_n(\omega)$. Since $\overline{X}_n(\omega)$ is a constant with respect to $\omega^* \in \Omega^*$, it follows that
$\mathbb{E}^*[X_{ni}(\omega)]=\overline{X}_n(\omega)-\overline{X}_n(\omega)=0.$ Similarly, $\mathbb{E}^*[||X_{n1}(\omega)||^2]=n^{-1} \sum_{i=1}^n ||X_i(\omega)-\overline{X}_n(\omega)||^2$, which as a finite sum of real numbers is finite. In the sequel we use the following notations:
If $x \in \mathbb{R}^d$ is a (column-) vector, then $x^{(k)}$ denotes its $k$-th component. If $M$ is a matrix, then $M_{kl}$ is its $(k,l)$-th entry. It follows that
\begin{eqnarray*}
& &(\text{Cov}^*[X_{n1}(\omega)])_{kl}= \text{Cov}^*[X_{n1}(\omega)^{(k)},X_{n1}(\omega)^{(l)}]=\text{Cov}^*[\xi_{n1}(\omega)^{(k)},\xi_{n1}(\omega)^{(l)}]\\
&=& E^*[\xi_{n1}(\omega)^{(k)} \xi_{n1}(\omega)^{(l)}]- E^*[\xi_{n1}(\omega)^{(k)}] E^*[\xi_{n1}(\omega)^{(l)}]\\
&=& n^{-1} \sum_{i=1}^n X_i(\omega)^{(k)}X_i(\omega)^{(l)}-\overline{X}_n(\omega)^{(k)}\overline{X}_n(\omega)^{(l)}
\rightarrow \Gamma_{kl} \quad \text{for } \mathbb{P}-\text{almost all } \omega \in \Omega
\end{eqnarray*}
by the Strong Law of Large Numbers (SLLN). Finally,
\begin{eqnarray} \label{Lindeberg}
0 &\le& L_n(\epsilon,\omega)= \mathbb{E}^*[||X_{n1}(\omega)||^2 1_{\{||X_{n1}(\omega)||> \epsilon \sqrt{n}\}}] \nonumber\\
&=& n^{-1}\sum_{i=1}^n ||X_i(\omega)-\overline{X}_n(\omega)||^2 1_{\{||X_i(\omega)-\overline{X}_n(\omega)||> \epsilon \sqrt{n}\}} \label{LB}\\
&\le& 2 n^{-1} \sum_{i=1}^n ||X_i(\omega)||^2 1_{\{||X_i(\omega)-\overline{X}_n(\omega)||> \epsilon \sqrt{n}\}}+2 ||\overline{X}_n(\omega)||^2, \label{LBbound}
\end{eqnarray}
where the last inequality follows from $||X_i(\omega)-\overline{X}_n(\omega)||^2 \le 2 (||X_i||^2+||\overline{X}_n(\omega)||^2)$. Since
$||\overline{X}_n(\omega)||^2 \rightarrow 0$ for $\mathbb{P}-\text{almost all } \omega \in \Omega$ by the SLLN and continuity of the norm, it remains to consider the first summand in (\ref{LBbound}) without the factor $2$, i.e.
$$
 \overline{L}_n(\epsilon,\omega):=n^{-1} \sum_{i=1}^n ||X_i(\omega)||^2 1_{\{||X_i(\omega)-\overline{X}_n(\omega)||> \epsilon \sqrt{n}\}} \ge 0.
$$
For that purpose notice that for every $m \in \mathbb{N}$ there exists a natural number $n_0=n_0(m)$ such that $\epsilon \sqrt{n} \ge m$ for all
$n \ge n_0$ and therefore
$$
 \overline{L}_n(\epsilon,\omega) \le n^{-1} \sum_{i=1}^n ||X_i(\omega)||^2 1_{\{||X_i(\omega)-\overline{X}_n(\omega)||> m\}} \quad \forall \; n \ge n_0 \; \forall\; m \in \mathbb{N} \; \forall \; \omega \in \Omega.
$$
Thus, if
$$
\Omega_0:= \liminf_{n \rightarrow \infty} \{||\overline{X}_n|| \le 1/2\}= \{||\overline{X}_n|| \le 1/2 \text{ for eventually all } n \in \mathbb{N}\},
$$
then $\mathbb{P}(\Omega_0)=1$ by the SLLN. (Here and in the sequel we omit the argument $\omega \in \Omega$.) Moreover, by $||X_i-\overline{X}_n||\le ||X_i||+ ||\overline{X}_n||$,
\begin{equation} \label{superset}
 \Omega_0 \subseteq \liminf_{n \rightarrow \infty} \{\overline{L}_n(\epsilon) \le n^{-1} \sum_{i=1}^n ||X_i||^2 1_{\{||X_i||> m/2\}}\}.
\end{equation}
Observe that the inclusion in (\ref{superset}) holds for all $m \in \mathbb{N}$, whence
$$
 \Omega_0 \subseteq \bigcap _{m \in \mathbb{N}} \liminf_{n \rightarrow \infty} \{\overline{L}_n(\epsilon) \le n^{-1} \sum_{i=1}^n ||X_i||^2 1_{\{||X_i||> m/2\}}\}.
$$
Consequently,
$$
 \Omega_0 \subseteq \{\limsup_{n \rightarrow \infty}\overline{L}_n(\epsilon) \le \limsup_{n \rightarrow \infty}n^{-1} \sum_{i=1}^n ||X_i||^2 1_{\{||X_i||> m/2\}} \; \forall \; m \in \mathbb{N}\}.
$$
Put
$$
  \Omega_m:=\{n^{-1} \sum_{i=1}^n ||X_i||^2 1_{\{||X_i||> m/2\}} \rightarrow \mathbb{E}[||X_1||^2 1_{\{||X_1||> m/2\}}], n \rightarrow \infty\}.
$$
Another application of the SLLN yields that $\Omega_\infty := \bigcap_{m \in \mathbb{N}} \Omega_m$ has probability $1$.
Let
$$
  E(m):= \mathbb{E}[||X_1||^2 1_{\{||X_1||> m/2\}}].
$$
Then
\begin{equation} \label{superset2}
 \Omega_0 \cap \Omega_\infty \subseteq \{ 0 \le \limsup_{n \rightarrow \infty} \overline{L}_n(\epsilon) \le E(m) \; \forall \; m \in \mathbb{N}\}.
\end{equation}
The Dominated Convergence Theorem ensures that $E(m) \rightarrow 0$ as $m \rightarrow \infty$. Thus taking the limit $m \rightarrow \infty$ in the event on the right side of (\ref{superset2}) shows that
$$
\{ 0 \le \limsup_{n \rightarrow \infty} \overline{L}_n(\epsilon) \le E(m) \; \forall \; m \in \mathbb{N}\} \subseteq \{\overline{L}_n(\epsilon) \rightarrow 0\}.
$$
Deduce from (\ref{superset2}) that $\overline{L}_n(\epsilon) \rightarrow 0 \; \mathbb{P}-$almost surely (a.s.), which in view of (\ref{LBbound}) results in
$L_n(\epsilon) \rightarrow 0$ a.s. for all $0<\epsilon \in \mathbb{Q}.$ Since $\mathbb{Q}$ is countable, we actually have
that $L_n(\epsilon) \rightarrow 0$ for all $0<\epsilon \in \mathbb{Q}$ a.s.
Now, Lemma \ref{convVn} yields the desired result (\ref{bootstrapconv}).
\end{proof}

\begin{equation} \label{arisesequence}
 X_{ni}=(Z_i^{(1)},\ldots,Z_i^{(d_n)}), 1 \le i \le n, \; n \in \mathbb{N}.
\end{equation}

\begin{equation} \label{norm}
 ||x||_{d_n}=||r_n(x)|| \quad \forall \; x \in \mathbb{R}^{d_n}.
\end{equation}

\begin{proposition} \label{CLTl2arrays} For every $n \in \mathbb{N}$ let $Z_{n1},\ldots,Z_{nn}$ i.i.d. random variables in $l^2$.
Assume that
\begin{itemize}
\item[(1)] $\mathbb{E}[Z_{n1}^{(i)}]=0$ and $\mathbb{E}[(Z_{n1}^{(i)})^2] < \infty \quad \forall \; i \in \mathbb{N} \; \forall \; n \in \mathbb{N}.$\\

\item[(2)] $\Gamma_n(k,l):= \mathbb{E}[Z_{n1}^{(k)} Z_{n1}^{(l)}] \rightarrow \Gamma(k,l) \in \mathbb{R}, \; n \rightarrow \infty \quad \forall \; k,l \in \mathbb{N}.$\\

\item[(3)] $\limsup_{n \rightarrow \infty} \sum_{k \ge 1} \Gamma_n(k,k) \le \sum_{k \ge 1} \Gamma(k,k) < \infty.$\\

\item[(4)] $\mathbb{E}[||\gamma_l(Z_{n1})||_l^2 1_{\{||\gamma_l(Z_{n1})||_l > \epsilon \sqrt{n}\}}] \rightarrow 0, \; n \rightarrow \infty \quad \text{for all rational } \epsilon>0.$\\
\end{itemize}
Then
\begin{equation} \label{CLTl2}
  \frac{1}{\sqrt{n}} \sum_{i=1}^n Z_{ni} \stackrel{\mathcal{D}}{\rightarrow} N(0,\Gamma) \quad \text{in } l^2,
\end{equation}
where $\Gamma= (\Gamma(k,l))_{k,l \in \mathbb{N}}.$
\end{proposition}

\begin{corollary} \label{normZnquer} Let $\overline{Z}_n := n^{-1} \sum_{i=1}^n Z_{ni}$ and $W_n:=n||\overline{Z}_n||^2.$ Then under the assumption of the above Proposition \ref{CLTl2arrays},
\begin{equation} \label{Vnpn}
 \sup_{x \in \mathbb{R}}|\mathbb{P}(W_n \le x - F_\infty(x)| \rightarrow 0, \; n \rightarrow \infty.
\end{equation}
Here, $F_\infty$ is the distribution function of $V_\infty = \sum_{i=1}^\infty \lambda_i N_i^2$, where $\lambda_1,\lambda_2,\ldots$ are the eigenvalues of $\Gamma$ and the $N_i, i \in \mathbb{N},$ are i.i.d. standard normal.
\end{corollary}

\begin{theorem} \label{bootstraphigh} Suppose $Z_i, i \in \mathbb{N},$ are i.i.d. random variables with values in $l^2$ such that
$\mathbb{E}[Z_1]=0$ and  $\mathbb{E}[||Z_1||^2] < \infty$. If $\Gamma = \text{Cov}[Z_1]$, then
\begin{equation} \label{uniformWnstar}
  \sup_{x \in \mathbb{R}}|\mathbb{P}^*(W_n^*(\omega) \le x)-F_\infty(x)| \rightarrow 0, \;  n \rightarrow \infty,\quad \text{for } \mathbb{P}\text{-almost all } \omega \in \Omega.
\end{equation}
Here, the $\lambda_i, i \in \mathbb{N}$ are the eigenvalues of $\Gamma$ and $N_i, i \in \mathbb{N}$ are i.i.d. standard normal. Moreover,  $F_\infty$ is the distribution function of $V_\infty=\sum_{i=1}^\infty \lambda_i N_i^2.$
\end{theorem}

\begin{proof} For each $n \in \mathbb{N}$ let $$\zeta_{nk}(\omega)(\omega^*):=Z_{i_{nk}(\omega^*)}(\omega), \; 1 \le k \le n$$ and
$$
\overline{\zeta}_n(\omega):= \frac{1}{n} \sum_{k=1}^n (\zeta_{nk}(\omega)- \overline{Z}_n(\omega))
$$
with $\overline{Z}_n(\omega):=n^{-1} \sum_{i=1}^n Z_i(\omega).$ Then by linearity and (\ref{arisesequence}), $\overline{\xi}_n(\omega) = c_n(\overline{\zeta}_n(\omega))$ and
therefore $||\overline{\xi}_n(\omega)||_{d_n}=||c_n(\overline{\zeta}_n(\omega))||_{d_n}$, which by (\ref{norm}) results in
$$
 W_n^*(\omega)=n||p_n(\overline{\zeta}_n(\omega))||^2.
$$
Since $p_n(\zeta_n(\omega))=n^{-1} \sum_{i=1}^n p_n(\zeta_{ni}(\omega)-\overline{Z}_n(\omega))$, we want to apply Corollary \ref{normZnquer} to
the array $Z_{ni}(\omega):= p_n(\zeta_{ni}(\omega)-\overline{Z}_n(\omega)), 1 \le i \le n, n \in \mathbb{N}.$ For that purpose notice that
$\zeta_{n1}(\omega),\ldots,\zeta_{nn}(\omega)$ are i.i.d. under $\mathbb{P}^*$ with common distribution
$$
 Q_n(\omega)=\frac{1}{n} \sum_{i=1}^n \delta_{Z_i(\omega)},
$$
where $\delta_x$ is the Dirac-measure at point $x \in l^2.$ Thus, $Z_{n1}(\omega),\ldots,Z_{nn}(\omega)$ are i.i.d. under $\mathbb{P}^*$ for every $n \in \mathbb{N}.$ Furthermore,
$$
 \mathbb{E}^*[Z_{n1}(\omega)]=\mathbb{E}^*[p_n(\zeta_{n1}(\omega)-\overline{Z}_n(\omega))]
 =p_n(\mathbb{E}^*[\zeta_{n1}(\omega)-\overline{Z}_n(\omega)])=p_n(0)=0,
$$
because by there definitions the operators $\mathbb{E}^*$ and $p_n$ can be commuted and $\mathbb{E}^*[\zeta_{n1}(\omega)]= \overline{Z}_n(\omega).$
In particular
\begin{equation} \label{Zn1arecentered}
\mathbb{E}^*[Z_{n1}(\omega)^{(i)}] = 0 \quad \forall \; i,n \in \mathbb{N} \; \forall \; \omega \in \Omega.
\end{equation}
In the sequel notice that by definition
\begin{equation} \label{zerocomponents}
   Z_{n1}(\omega)^{(i)} = \zeta_{n1}(\omega)^{(i)}-\overline{Z}_n(\omega)^{(i)} \quad \forall \; i \le d_n \quad \text{and} \quad                           Z_{n1}(\omega)^{(i)} =0 \quad \forall \; i >d_n.
\end{equation}
So, if $i \le d_n$, then
$$
 \mathbb{E}^*[(Z_{n1}(\omega)^{(i)})^2]=\int_{l^2} (x^{(i)}-\overline{Z}_n(\omega)^{(i)})^2 Q_n(\omega)(dx)= n^{-1} \sum_{j=1 }^n (Z_j(\omega)^{(i)}-\overline{Z}_n(\omega)^{(i)})^2,
$$
i.e. the second moment is equal to the sample variance. With a view to (\ref{zerocomponents}) we obtain that
\begin{equation} \label{Zn1secondmoments}
\mathbb{E}^*[(Z_{n1}(\omega)^{(i)})^2] < \infty \quad \forall \; i,n \in \mathbb{N} \; \forall \; \omega \in \Omega.
\end{equation}
Similarly, for every fixed natural numbers $k,l \le d_n$ it follows that
\begin{eqnarray*}
 \Gamma_n(\omega)(k,l)&:=& \text{Cov}^*[Z_{n1}(\omega)^{(k)}, Z_{n1}(\omega)^{(l)}]=\text{Cov}^*[\zeta_{n1}(\omega)^{(k)}-\overline{Z}_n(\omega)^{(k)},\zeta_{n1}(\omega)^{(l)}-\overline{Z}_n(\omega)^{(l)}]\nonumber\\
 &=&\text{Cov}^*[\zeta_{n1}(\omega)^{(k)},\zeta_{n1}(\omega)^{(l)}]\nonumber\\
 &=& n^{-1} \sum_{i=1}^n Z_i(\omega)^{(k)} Z_i(\omega)^{(l)}- \overline{Z}_n(\omega)^{(k)} \overline{Z}_n(\omega)^{(l)}.
\end{eqnarray*}
Since $k,l \le d_n$ for eventually all $n \in \mathbb{N}$ as $d_n \rightarrow \infty$, the SLLN yields that
\begin{equation} \label{convergenceofthecovariances}
 \Gamma_n(\omega)(k,l) \rightarrow (\text{Cov}[Z_1])_{k,l} \quad \forall \; k,l \in \mathbb{N} \quad \text{for }\mathbb{P}-\text{almost all } \omega \in \Omega.
\end{equation}

Observe that by the Dominated Convergence Theorem,
\begin{eqnarray}
 \sum_{k \ge 1}\Gamma_n(\omega)(k,k)&=& \sum_{k \ge 1}\mathbb{E}^*[Z_{n1}(\omega)^2]= \mathbb{E}^*[||Z_{n1}(\omega)||^2]=\mathbb{E}^*[||p_n(\zeta_{n1}(\omega))-p_n(\overline{Z}_n(\omega)||^2] \nonumber\\
 &=& n^{-1}\sum_{i=1}^n ||p_n(Z_i(\omega))-p_n(\overline{Z}_n(\omega)||^2. \label{sumnormsquared}
\end{eqnarray}
Let $<\cdot,\cdot>$ denote the scalar product in $l^2$. Then
$$
||p_n(Z_i(\omega))-p_n(\overline{Z}_n(\omega)||^2 = ||p_n(Z_i(\omega))||^2 - 2<p_n(Z_i(\omega)),p_n(\overline{Z}_n(\omega)>+||p_n(\overline{Z}_n(\omega)||^2
$$
Thus (\ref{sumnormsquared}) and linearity of the scalar product yield
\begin{equation} \label{upperboundsumGamman}
 \sum_{k \ge 1}\Gamma_n(\omega)(k,k)= n^{-1}\sum_{i=1}^n ||p_n(Z_i(\omega))||^2-||p_n(\overline{Z}_n(\omega)||^2 \le  n^{-1}\sum_{i=1}^n ||Z_i(\omega))||^2,
\end{equation}
because $||p_n(x)|| \le ||x||$ for all $x \in l^2.$ By the SLLN the upper bound in (\ref{upperboundsumGamman}) converges to $\mathbb{E}[||Z_1||^2]$ for $\mathbb{P}-$almost every $\omega \in \Omega$. Since $\mathbb{E}[||Z_1||^2]=\sum_{k \ge 1}\Gamma(k,k)$ by the Monotone Convergence Theorem, we therefore obtain
\begin{equation} \label{tracecondition}
 \limsup_{n \rightarrow \infty}\sum_{k \ge 1}\Gamma_n(\omega)(k,k) \le \sum_{k \ge 1}\Gamma(k,k)<\infty \quad \text{for }\mathbb{P-}\text{almost all } \omega \in \Omega,
\end{equation}
where finiteness of the series follows from $\mathbb{E}[||Z_1||^2]< \infty$ by assumption. Finally, notice that
$\gamma_l(Z_{n1}(\omega))= \gamma_l(\zeta_{n1}(\omega))-\gamma_l(\overline{Z}_n(\omega)) \in \mathbb{R}^l$ for eventually all $n \in \mathbb{N}$.
For all these $n$ the Lindeberg-term is equal to
\begin{eqnarray*}
 L_n(\epsilon,\omega)&=&\mathbb{E}^*[||\gamma_l(Z_{n1}(\omega))||^2 1_{\{||\gamma_l(Z_{n1}(\omega))||> \sqrt{n} \epsilon\}}\\
  &=& n^{-1}\sum_{i=1}^n ||\gamma_l(Z_i(\omega))-\gamma_l(\overline{Z}_n(\omega))||^2 1_{||\gamma_l(Z_i(\omega))-\gamma_l(\overline{Z}_n(\omega))||>\sqrt{n} \epsilon\}}.
\end{eqnarray*}
This is equation (\ref{LB}) with $X_i$ there corresponds to $\gamma_l(Z_i)$ here. From this point (\ref{LB}) on, we can use the same line of reasoning to obtain
\begin{equation} \label{LBcondition}
 L_n(\epsilon,\omega) \rightarrow 0 \quad \text{for all rational } \epsilon >0 \quad \text{for }\mathbb{P}\text{-almost all } \omega \in \Omega.
\end{equation}

With (\ref{Zn1arecentered}), (\ref{Zn1secondmoments}), (\ref{convergenceofthecovariances}), (\ref{tracecondition}) and (\ref{LBcondition})
all requirements of Corollary \ref{normZnquer} are fulfilled for $\mathbb{P}$-almost all $\omega \in \Omega$, which yields the desired result.
\end{proof}

\end{document}
